\documentclass[12pt,a4paper]{article}
% Preambule commun aux comptes rendus CR_*.tex.

\usepackage{array}
\newcolumntype{C}[1]{>{\centering\arraybackslash}m{#1}}
\usepackage{multirow}
\usepackage{geometry}
\geometry{a4paper, margin=1in}
\usepackage{graphicx}
\usepackage{float}
\usepackage{tikz}
\usetikzlibrary{arrows.meta,positioning}
\usepackage{pgfgantt}
\usepackage{pdflscape}
\usepackage{enumitem}
\usepackage{booktabs}
\usepackage{longtable}
\usepackage{ragged2e}
\usepackage{tabularx}
\usepackage{xparse}
\usepackage{ltablex}
\usepackage{xltabular}
\usepackage{subcaption}

\newcolumntype{Y}{>{\RaggedRight\arraybackslash}X}

\newcommand{\StyledTableSetup}{%
  \footnotesize
  \renewcommand{\arraystretch}{1.2}%
  \setlength{\tabcolsep}{4pt}%
}


% ============================================================
% IHEDN COVER TEMPLATE (XeLaTeX)
% ============================================================
% Compilation (from project root):
%   latexmk -pdf main.tex
%   LATEXMK_SHELL_ESCAPE=0 latexmk -pdf main.tex
%
% Optional environment controls from latexmkrc:
%   LATEXMK_ENGINE=xelatex        (default/enforced)
%   LATEXMK_SHELL_ESCAPE=1|0      (default 1, needed for SVG conversion)
%
% Required package for this template:
%   \usepackage{ihedn-cover}
%
% Note: ihedn-cover already loads needed internal packages
% (fontspec, graphicx, tikz, ragged2e, optional svg).
% ============================================================

% ----------------------------------------------------------------
% Package options (documented)
% ----------------------------------------------------------------
% Geometry scope:
%   apply-geometry            -> force A4 + zero margins while cover is built
%   no-apply-geometry         -> do not alter host geometry (default)
%
% Typesetting freeze:
%   global-typesetting        -> apply freeze globally (intrusive)
%   no-global-typesetting     -> no global freeze (default)
%
% Small-caps handling:
%   local-smallcaps-fix       -> \textsc -> \MakeUppercase only inside cover (default)
%   no-local-smallcaps-fix    -> disable local fix
%   global-smallcaps-fix      -> global override (intrusive)
%   no-global-smallcaps-fix   -> disable global fix (default)
%
% Main font handling:
%   global-mainfont           -> set Times New Roman globally (default)
%   local-mainfont            -> set Times only during cover rendering
%   no-mainfont-override      -> keep host document main font
%
% SVG support:
%   svg-support               -> enable SVG logos (default, needs shell-escape)
%   no-svg-support            -> disable SVG path handling (pdf/png only)
%
% Debug overlay:
%   debug                     -> show mm grid + bounding boxes on cover
%   no-debug                  -> hide debug overlay (default)
% ----------------------------------------------------------------
\usepackage[
  no-apply-geometry,
  no-global-typesetting,
  global-smallcaps-fix,
  local-smallcaps-fix,
  global-mainfont,
  svg-support,
  no-debug
]{ihedn-cover}

% ----------------------------------------------------------------
% Bibliography stack (BibLaTeX/Biber, style from subtree vendor/)
% ----------------------------------------------------------------
\usepackage{polyglossia}
\setdefaultlanguage{french}
\makeatletter
\g@addto@macro\captionsfrench{%
  \renewcommand{\tablename}{Tableau}%
  \renewcommand{\figurename}{Figure}%
}
\makeatother
\usepackage{csquotes}
\usepackage{xurl}
\usepackage{hyperref}
\makeatletter
\let\ihedn@orig@input@path\input@path
\def\input@path{{vendor/biblatex-sciences-po-fr/style/}}
\makeatother
\usepackage[
  backend=biber,
  style=sciencespo,
  hyperref=true,
  autolang=other
]{biblatex}
\AtBeginBibliography{\sloppy\setlength{\emergencystretch}{3em}}
\makeatletter
\let\input@path\ihedn@orig@input@path
\makeatother
\addbibresource{bibliographies/bib/rapport.bib}

% ----------------------------------------------------------------
% tcolorbox
% ----------------------------------------------------------------
\usepackage[most]{tcolorbox}
\tcbuselibrary{breakable}

% Example: use Arial for the full report + cover.
% Uncomment and switch package option to no-mainfont-override.
% \setmainfont{Arial}[
%   UprightFont    = *,
%   BoldFont       = * Bold,
%   ItalicFont     = * Italic,
%   BoldItalicFont = * Bold Italic
% ]

%==================================================
% Liste des propositions
%==================================================

\makeatletter

\newcommand{\listpropositionname}{Liste des propositions}

\newcommand{\listofpropositions}{
  \section*{\listpropositionname}
  \addcontentsline{toc}{section}{\listpropositionname}
  \@starttoc{lop}
}

\newcommand*\l@proposition{\@dottedtocline{1}{1.5em}{2.3em}}

\makeatother

%==================================================
% Compteur des propositions
%==================================================

\newcounter{proposition}

%==================================================
% Style graphique des propositions
%==================================================

\newtcolorbox{propositionbox}[2][]{
  lower separated=false,
  colback=white!80!gray,
  colframe=white!20!black,
  fonttitle=\bfseries,
  colbacktitle=white!30!gray,
  coltitle=black,
  enhanced,
  sharp corners,
  attach boxed title to top left={
    xshift=0.5cm,
    yshift=-2mm
  },
  title=#2,
  #1
}

%==================================================
% Environnement proposition
%==================================================

\newenvironment{proposition}[1]
{
  \refstepcounter{proposition}

  \addcontentsline{lop}{proposition}{%
    \protect\numberline{\theproposition}#1%
  }

  \begin{propositionbox}
  {Proposition~\theproposition~: #1}
}
{
  \end{propositionbox}
}

% Renommage Liste des figures en Table des figures
\addto\captionsfrench{%
  \renewcommand{\listfigurename}{Table des figures}
}

\newcommand{\heure}[2]{{\NoAutoSpacing #1:#2}}

\DeclareFieldFormat[misc]{title}{#1}

% Activer ce package avec l'option none retire la césure.
%\usepackage[none]{hyphenat}

\usepackage{adjustbox}
\usepackage{pdfpages}

\makeatletter
\newcommand{\IhednSetPdfPageCount}[2]{%
  \begingroup
  \edef\ihedn@pdfpath{#2}%
  \ifXeTeX
    \global#1=\XeTeXpdfpagecount "\ihedn@pdfpath" %
  \else
    \ifLuaTeX
      \saveimageresource{\ihedn@pdfpath}%
      \global#1=\lastsavedimageresourcepages\relax
    \else
      \pdfximage{\ihedn@pdfpath}%
      \global#1=\pdflastximagepages\relax
    \fi
  \fi
  \endgroup
}
\makeatother

\newcount\IhednIncludedPdfPageCount
\newcommand{\IhednRenderPdfPagesInBoxes}[1]{%
  \IhednSetPdfPageCount{\IhednIncludedPdfPageCount}{#1}%
  \foreach \IhednIncludedPdfPage in {1,...,\the\IhednIncludedPdfPageCount}{%
    \begin{tcolorbox}[
      colback=gray!5,
      colframe=black,
      boxrule=0.5pt,
      arc=3mm
    ]
    \centering
    \includegraphics[page=\IhednIncludedPdfPage,width=0.95\textwidth]{#1}
    \end{tcolorbox}
    \ifnum\IhednIncludedPdfPage<\IhednIncludedPdfPageCount
      \clearpage
    \fi
  }%
}

\usepackage{tablefootnote}


\begin{document}

% ----------------------------------------------------------------
% Cover keys (all available keys are shown below)
% ----------------------------------------------------------------
% Header block (top-right), logo, title, report lines, subtitle, committee.
\PageDeGardeIHEDN{
    header_1={Union-IHEDN},
    header_2={Association des auditeurs IHEDN},
    header_3={région Paris / Île de France},
    date={le \today},
    logo_path={Figures/logos_IHEDN/Logo-AR16},
    title={Crises majeures : quel rôle pour les citoyens engagés et les associations IHEDN ?},
    title_max_lines=2,
    title_fontsize={26.000pt},
    title_baselineskip={28.667pt},
    reportline_1={Rapport du Comité d'étude de l'Association régionale des auditeurs},
    reportline_2={de l'IHEDN région Paris / Île de France (AR 16)},
    subtitle={Présentation de la problématique et du cadre général du questionnaire en vue d'établir une cartographie claire des membres de l'AR 16.},
    committee={Comité d'étude, d'octobre 2025 à juin 2026, composé de : Mme Valérie Arlé-Faivre, \textit{présidente} ; M. Aurélien Duvignac-Rosa, \textit{rapporteur} ; M. Pascal Bayot Duponchel ; M. Jean-Jacques Cleynen ; Mme Valérie Cowan ; M. Alain Fontan ; M. Pierre-Marie Giard ; Mme Muriel Joyeux ; M. Bruno Leray ; M. Yves-Henri Renhas ; M. Marc Roure ; Mme Isabelle de Segonzac ; M. Gérard Turck ; Mme Marie Danielle Vasquez-Duchêne.}
}

\clearpage
\pagenumbering{Roman}
\tableofcontents
\clearpage
\pagenumbering{arabic}

\section{Objectifs de l'enquête ?}

À l'occasion du 50\textsuperscript{e} anniversaire de l'AR 16\footcite{linkedin_AR16}, la question de son positionnement stratégique et de sa capacité d'engagement dans la gestion des crises majeures se pose avec une acuité particulière. 

Or, toute réflexion stratégique suppose un préalable : disposer d'une connaissance objectivée de l'Association et de ses membres. À ce jour, notre perception des ressources, des compétences et des disponibilités repose largement sur des représentations partielles.

L'enquête proposée vise à produire un état des lieux fiable permettant de dépasser ces impressions et de fonder les décisions futures sur des données empiriques.

\section{Problématique centrale}

La question directrice est la suivante :

\begin{quote}
\textbf{Quelle est la capacité réelle et mobilisable de l'AR 16 à contribuer à la préparation, à l'accompagnement et à l'analyse des crises majeures ?}
\end{quote}

L'hypothèse sous-jacente est que l'Association dispose d'un capital de compétences significatif, mais insuffisamment cartographié et structuré. Le questionnaire constitue l'outil permettant d'objectiver cette hypothèse.

\section{Finalités stratégiques}

L'enquête poursuit quatre objectifs :

\begin{itemize}
    \item Cartographier précisément les profils et compétences des membres ;
    \item Identifier les conditions réelles d'engagement (contraintes professionnelles, disponibilité, statut) ;
    \item Renforcer la crédibilité institutionnelle de l'AR 16 par des données objectivées ;
    \item Structurer l'Association (mise à jour de l'annuaire, identification d'expertises spécifiques).
\end{itemize}

\newpage

\section{Cadre méthodologique}

Le questionnaire devra découler d'une problématique explicite et d'hypothèses clairement formulées. Chaque question devra répondre à un objectif analytique précis.

Deux précautions sont essentielles :

\begin{itemize}
    \item Éviter toute \enquote{imposition de problématique} susceptible d'orienter les réponses\footcite{Bourdieu1980} \footcite[70]{DeSingly2020} \footcite[84]{Fenneteau2015} ;
    \item Interpréter les données biographiques avec prudence, en tenant compte des déterminations sociales et professionnelles\footcite[45]{Bourdieu1980}.
\end{itemize}

Le questionnaire sera structuré de manière progressive, précédé d'un paragraphe explicatif, et d'une longueur raisonnable afin de garantir un taux de réponse satisfaisant.

\section{Cadre juridique}

La mise en œuvre suppose :

\begin{itemize}
    \item la conformité aux exigences du RGPD ;
    \item la désignation d'un référent pour la protection des données ;
    \item la rédaction d'un formulaire de consentement ;
    \item l'anonymisation des réponses.
\end{itemize}

\section{Enjeu pour le CODIR}

\textbf{Valider le principe de cette enquête revient à choisir de fonder le positionnement stratégique de l'AR 16 sur une connaissance objectivée plutôt que sur des représentations partielles.}

Le questionnaire constitue ainsi un \textbf{levier de structuration interne}, \textbf{de clarification stratégique} et \textbf{de valorisation institutionnelle}.

\bigskip

\noindent
\textbf{Il ne s'agit pas d'un simple outil administratif, mais d'un instrument d'analyse et de projection stratégique.}

\newpage

\printbibliography[heading=bibintoc,title={Références bibliographiques}]

% ----------------------------------------------------------------
% Notes d'usage:
% - Le style bibliographique est dans:
%   vendor/biblatex-sciences-po-fr/style/
% - Ajouter des entrées dans les fichiers bibliographiques chargés ci-dessus:
%   vendor/biblatex-sciences-po-fr/bibliographies/{articles,livres,online,rapports,theses}.bib
% - Compilation requise:
%   latexmk -pdf main.tex
%   (latexmkrc est configuré pour XeLaTeX + Biber).
% ----------------------------------------------------------------

\end{document}
